\section{总结与延伸}

Week 9 将视角从``写代码''转向``运行代码''。DevOps 和 SRE 是软件工程中最耗时的部分，AI SRE 正在通过知识图谱、跨栈 Agent 和实时分析来降低运维负担。

\subsection{关键要点}

\begin{itemize}
    \item 编码只占工程时间的 30\%，运维是更难的 70\%
    \item 四个黄金信号：延迟、错误、流量、饱和度
    \item 事件响应遵循标准化 playbook，核心度量是 MTTR
    \item AI SRE 的四大特征：动态知识图谱、跨栈 Agent、实时叙述、可解释性
    \item 最大价值：打破知识孤岛，降低 MTTR
    \item 当前从根因分析入手，自动修复仍在早期
\end{itemize}

\subsection{运维落地清单}

\begin{enumerate}
    \item 先统一监控口径和日志质量，再引入 AI SRE，不要让 AI 替代混乱的数据基础
    \item 把 MTTR、误报率和升级准确率作为核心指标，而不是只看模型回答是否“像样”
    \item 在受控范围内自动化低风险操作，把高风险修复留给人工审批
\end{enumerate}

\subsection{值班角色分工表}

\begin{table}[H]
\centering
\begin{tabular}{p{0.18\textwidth} p{0.26\textwidth} p{0.3\textwidth}}
\toprule
角色 & 在 AI SRE 体系中的职责 & 不能放弃的人类责任 \\
\midrule
值班工程师 & 审阅 AI 汇总、决定缓解动作、协调升级 & 对重大操作负责最终确认 \\
AI SRE 系统 & 聚合证据、给出根因假设、建议低风险动作 & 不能越权执行高风险修复 \\
平台/基础设施团队 & 维护监控、日志、runbook 与权限边界 & 保证 AI 有高质量运行数据可用 \\
\bottomrule
\end{tabular}
\caption{AI SRE 场景下的人机角色边界}
\end{table}

\subsection{Postmortem 闭环为什么更重要了}

AI SRE 并不会让事故复盘变得不重要，反而会让它更重要。因为每一次事故处理，都会同时暴露两类问题：

\begin{itemize}
    \item 系统本身的问题：监控缺口、依赖不清、runbook 过时、权限边界模糊
    \item AI 系统的问题：它遗漏了哪些信号、放大了哪些噪音、在哪一步给出了不可靠建议
\end{itemize}

因此，成熟团队需要把 postmortem 的产出重新写回监控、告警规则、AI 提示词、runbook 和审批边界，让 AI SRE 不是一次性的事故助手，而是可持续进化的运维系统。

\subsection{拓展阅读}

\begin{itemize}
    \item Resolve AI：\url{https://resolve.ai}
    \item Google SRE Book：\url{https://sre.google/sre-book/table-of-contents/}
    \item DataDog Bits AI：\url{https://www.datadoghq.com/product/platform/bits-ai/}
    \item ``The Four Golden Signals'' (Google SRE Book, Chapter 6)
    \item PagerDuty Incident Response Guide：\url{https://response.pagerduty.com}
\end{itemize}

%% --- Fin del contenido --- %%

\end{document}
